% ============================================================
% ICS & AI Transparency Log
% Compile: tectonic -X compile ICS.tex
% ============================================================
\documentclass[11pt]{article}
\usepackage[margin=2.3cm]{geometry}
\usepackage{amsmath,amssymb}
\usepackage{booktabs}
\usepackage{fontspec}
\usepackage{microtype}
\usepackage{xcolor}
\usepackage[hidelinks]{hyperref}
\usepackage{enumitem}

\setmainfont{Times New Roman}
\newcommand{\e}[1]{{\emojifont #1}}

\setlength{\parindent}{0pt}
\setlength{\parskip}{0.5em}
\renewcommand{\arraystretch}{1.3}

\begin{document}

\begin{center}
  {\LARGE\bfseries ICS \& AI Transparency Log}\\[0.4em]
  {\large Appended to the end of the project report}
\end{center}
\vspace{0.2em}
\hrule
\vspace{0.6em}

Two documents in one file:
\begin{itemize}[leftmargin=2em, itemsep=1pt]
  \item \textbf{Part A --- Team AI Transparency Log} (part of rubric dimension E,
    Writing, Presentation \& AI Transparency, 15\%). One log per team, filled in
    collectively.
  \item \textbf{Part B --- Individual Contribution Statement (ICS)}. One statement per
    member, filled in individually.
\end{itemize}

% ============================================================
\section*{Part A --- Team AI Transparency Log}
% ============================================================
\textbf{Purpose.} AI tools are \textbf{encouraged} --- this log is not a
confession, it is a record of the gap between AI output and your understanding.
A team that used AI extensively but can explain every modification earns full
marks; a team that wrote everything by hand but cannot answer basic questions
does not.

\subsection*{A1. AI tools used}
\begin{table}[h]
\centering
\small
\begin{tabular}{@{}p{3.5cm}p{4.5cm}p{3.5cm}@{}}
\toprule
Tool & Version / Date & Primary use \\
\midrule
 & & \\
 & & \\
\bottomrule
\end{tabular}
\end{table}

\subsection*{A2. Significant AI-assisted claim audits}
Select two or three interactions that materially affected a mathematical claim,
algorithm, experiment, or interpretation. Routine autocomplete and spelling
corrections may be summarised once. Store raw transcripts separately; provide a
link or content hash if requested rather than appending pages of raw output.
\begin{table}[h]
\centering
\small
\begin{tabular}{@{}p{0.8cm}p{5.7cm}p{5.0cm}@{}}
\toprule
Item & Prompt / claim / assumptions & Risk and possible falsifier \\
\midrule
1 & & \\
2 & & \\
3 & & \\
\bottomrule
\end{tabular}
\end{table}

\subsection*{A3. Independent verification and decision}
\begin{table}[h]
\centering
\small
\begin{tabular}{@{}p{0.8cm}p{5.7cm}p{5.0cm}@{}}
\toprule
Item & Independent check and evidence location & Final decision / modification / limitation \\
\midrule
1 & & \\
2 & & \\
3 & & \\
\bottomrule
\end{tabular}
\end{table}
\emph{If you accepted an AI output unchanged, the independent check and its
evidence are still required. Quality of audit matters; number of prompts does not.}

\subsection*{A4. Intellectual ownership declaration}
Each team member signs below, identifying what the team owns and the trusted
dependencies it calls. Members affirm that they can explain and modify the
team-owned numerical methods, justify the mathematical claims, and state the
assumptions and tests at library interfaces without AI assistance during the
viva. The viva may ask for a small live modification or diagnostic test.
\begin{table}[h]
\centering
\small
\begin{tabular}{@{}p{2.4cm}p{2.4cm}p{4.0cm}p{3.2cm}@{}}
\toprule
Name & Signature & Team-owned artefact / contribution & Trusted dependency and interface \\
\midrule
 & & \\
 & & \\
 & & \\
\bottomrule
\end{tabular}
\end{table}

\subsection*{Example of a high-quality entry}
\begin{quote}
\textbf{Item 2: Adaptive step-size controller.}

\textbf{Prompt:} ``Write a Python function for adaptive RK4 step-size control
using embedded RK45 for the Van der Pol oscillator with $\mu = 1000$.''

\textbf{AI output:} a 40-line function with a simple error estimate.

\textbf{What we changed:} (1) replaced the fixed safety factor $S = 0.9$ with a
PI-controller formula based on Hairer--N{\o}rsett--Wanner (1993) \S II.4;
(2) added a stiffness detector that switches to implicit Euler after 3
consecutive rejected steps; (3) replaced the AI's (incorrect) global error
estimate with local error-per-step control.

\textbf{Why:} the AI controller is a textbook template for non-stiff problems;
for $\mu = 1000$ it took 42{,}000 steps vs.\ our 1{,}200 for $t \in [0, 50]$.
\end{quote}

% ============================================================
\section*{Part B --- Individual Contribution Statement (ICS)}
% ============================================================
One page per member, appended after the AI Transparency Log. Purpose:
accountability. The instructor reads these against the code (git history), the
report, and the viva answers.

\textbf{Not a team document.} Five different ICSs, five different voices. If two
ICSs claim the same paragraph of code, the instructor will ask both of you about it.

\subsection*{ICS format}
\begin{verbatim}
Individual Contribution Statement — Project 1 (ODE IVP)

Name:
Role this project:

The specific things I did this project:
1. [concrete item — e.g. ``Implemented the implicit Euler Newton solver in
   methods.py; verified quadratic convergence'']
2. [concrete item — e.g. ``Derived the stability region of RK4 in section 3.2
   of the report; wrote the finite-difference check'']
3. [concrete item — e.g. ``Wrote the log-log convergence script and the error
   analysis in section 5'']

What I would do differently:
[1-3 sentences]

Anything I should have asked for help with earlier:
[honest reflection, optional but valuable]
\end{verbatim}

\subsection*{Rules}
\begin{enumerate}[leftmargin=2.2em, itemsep=2pt]
  \item \textbf{Be specific.} ``I worked on the report'' is not a contribution.
    ``I wrote the stability section (3.2) including the derivation and the
    figure'' is.
  \item \textbf{Name files and functions.} The instructor will grep the git
    history --- your claims should be traceable to commits you made.
  \item \textbf{Be honest about the balance.} If you contributed less than your
    teammates, say so and say why. Honest self-assessment is worth more to the
    grade than inflated claims, which fail the viva.
  \item \textbf{AI use is already declared in Part A.} Cross-reference your ICS
    items to the Part A log entries where relevant.
\end{enumerate}

% ============================================================
\section*{Instructor's use}
% ============================================================
\begin{itemize}[leftmargin=2em, itemsep=2pt]
  \item \textbf{Part A:} assess the quality of the two or three claim audits:
    assumptions, risk, independent verification, evidence location, decision,
    and limitations. Do not reward prompt volume. A missing log caps rubric E
    at 2; a concise, rigorous log can earn full marks.
  \item \textbf{Part B:} cross-check claims against \texttt{git log} (who
    committed which file).
  \item In the final viva, ask each member to explain one thing \emph{outside}
    their own role and perform a small modification or diagnostic test on
    team-owned work. Ask what is delegated to libraries and how that interface
    was tested.
  \item The ICS, together with the viva and the issue-ticket record,
    determines the $\pm$ adjustment to the project grade (see assessment
    strategy in the teaching plan).
\end{itemize}

\end{document}
